\section{Why we cache, and what it weighs}
\label{sec:caches}

\subsection{The structure of redundant work}

A continuous-energy Monte Carlo case loads $\mathcal{O}(20)$
nuclides at $\mathcal{O}(10^5)$ energy points per nuclide; each
nuclide carries multiple reaction-channel SVD bases / pointwise
tables, an elastic angular distribution, fission outgoing-energy
distributions, optional URR probability tables, and (for
moderators) a separately-loaded S($\alpha$,$\beta$) thermal-scattering
library with $\mathcal{O}(10^2)$ incident energies, $\mathcal{O}(10^4)$
outgoing-energy bins, and $\mathcal{O}(10^5)$ angle bins per temperature.
The host parses these from HDF5; the device wants them in
flat-pack form ready for kernel pointer arithmetic. Both layers cost
real time.

An ICSBEP corpus sweep does not load $N$ distinct nuclides; it loads
the same $\sim$25 nuclides over and over. Every PWR pin cell uses
H-1, O-16, U-235, U-238, Zr-90 through Zr-96; every Be-reflected
fast-metal case uses Be-9 plus its \texttt{c\_Be} TSL; every HEU
solution case uses the same actinide block. A 375-case corpus at
5 seeds per case will, naively, parse \texttt{U235.h5} from disk
\textbf{1\,875 times} and HtoD-copy its SVD basis the same number
of times.

We measured this on the \texttt{origin/main} state immediately
before this paper's work landed. On a 73-case sub-sweep,
\texttt{c\_Be} was parsed from disk \textbf{125 times} and re-uploaded
to the GPU the same number of times --- roughly 8\,GB of redundant
HtoD copies for one TSL file alone.

\subsection{Host-side: \texttt{thermal\_cache} + \texttt{TieredStore}}

The fix on the host side is straightforward and has been in the
engine for a while:

\begin{itemize}
    \item \texttt{material\_resolve::thermal\_cache} is a
        process-wide \texttt{RwLock<HashMap<PathBuf,
        Arc<ThermalScatteringData>>>}. Keyed on the canonical
        absolute path of the TSL HDF5 file. First parse populates
        the entry; every subsequent \texttt{resolve\_materials} call
        for the same TSL hands back an \texttt{Arc::clone}. No
        size cap --- TSL inventory per data dir is bounded
        ($\leq$ 20 distinct files) and the bytes are immutable.
    \item \texttt{transport::nuclide\_cache::TieredStore} is the
        analogous structure for nuclide kernels, with an L1
        in-memory tier (\texttt{Arc<NuclideKernels>}-keyed) and an
        eviction policy combining LFU and recency
        (\texttt{LfuEntries::evict\_to\_budget}).
\end{itemize}

The host caches use \texttt{Arc::as\_ptr} (or canonical paths) as
identity keys. The \texttt{Arc::as\_ptr} approach matters because
the underlying \texttt{NuclideKernels} or
\texttt{ThermalScatteringData} structs are large
($\mathcal{O}(50)$-MB) and content-hashing them every cache lookup
would defeat the cache's purpose; pointer identity is a safe proxy
for content identity \emph{provided} every consumer goes through the
same upstream loader. The loader's job is to return the same
\texttt{Arc} for the same logical request --- a contract we enforce
by funnelling every consumer through \texttt{TieredStore::get\_or\_load}
or \texttt{thermal\_cache::load\_or\_get}.

\subsection{Device-side: \texttt{per\_nuclide\_cache} and \texttt{sab\_buffer\_cache}}

The device side needs a parallel cache structure because Arc
pointer identity does not give you the device buffer; it gives you
the host data the device buffer was once derived from. The first
call has to upload; subsequent calls have to find the cached upload
again. Our pattern, identical between the two caches:

\begin{enumerate}
    \item \textbf{Key} is the \texttt{Arc::as\_ptr} of the host
        \texttt{NuclideKernels} or \texttt{ThermalScatteringData},
        cast to \texttt{usize}, plus any structural parameters
        the kernel ABI bakes in: SVD rank for nuclides;
        \texttt{n\_nuc} (sizes the per-nuclide lookup table) plus
        the full slot list for SAB.
    \item \textbf{Value} is an \texttt{Arc<GpuNuclideData>} (or
        \texttt{Arc<GpuSabData>}). Returning an Arc means a cache
        hit is one atomic refcount bump --- no device traffic and
        no host-side rebuild of flat-pack arrays.
    \item \textbf{Eviction} is LFU-with-recency from
        \texttt{nuclide\_cache::eviction::evict\_to\_budget}.
        Both caches share the same \texttt{bundle\_cache\_budget}
        ($75\,\%$ of detected VRAM, env-var overridable as
        \texttt{OPEN\_RUST\_MC\_GPU\_BUNDLE\_CACHE\_BYTES} or
        \texttt{\_FRACTION}). The map / adapter is per-cache so an
        eviction in one does not perturb the other.
    \item \textbf{Race handling} is the usual ``check, compute
        outside the lock, then re-check on insert'' three-phase
        protocol. If two threads race the same key during a sweep
        preload, the loser drops its upload and clones the winner's
        \texttt{Arc}. Both buffers being byte-identical (same
        upstream Arc, same upload pipeline) means the dropped
        upload is wasted work, not wasted correctness.
\end{enumerate}

\subsection{Order-sensitivity in the SAB key}

S($\alpha$,$\beta$) takes one extra design decision the per-nuclide
cache did not need. The \texttt{upload\_sab\_data\_multi} routine
packs $N$ TSLs into a single flat layout, with per-slot offset
arrays (\texttt{slot\_inc\_e\_off},
\texttt{slot\_eout\_table\_off}, $\ldots$) that index into the
concatenated buffers. The slot order is the iteration order of the
input slice. Two callers asking for the same logical slots in
different orders produce different per-slot offset arrays, which
means different cached device buffers.

We make the cache key order-sensitive: it stores
\texttt{Vec<(arc\_ptr, nuc\_idx)>} in the upload order, not a
sorted set. Two callers can construct the slot list however they
like, but a sweep that consistently iterates
\texttt{provider.thermal} in slot-index order (which is what every
production caller does) gets clean cache hits. A future caller
that wants order-independence can sort upstream of the call site.

\subsection{What it weighs (measured)}

We ran the smoke-verification sweep on heu-comp-inter-003\_case-1
with 2 seeds, 30 batches, 5\,000 particles, ENDF/B-VIII.1
nuclear data, on the dev-box RTX~A1000 (4~GB), and snapshotted
the GPU caches before and after the case via the new
\texttt{-{}-debug-metrics} path (§\ref{sec:observability}).

\noindent\resizebox{\columnwidth}{!}{%
\begin{tabular}{@{}lrr@{}}
\toprule
phase                            & before     & after      \\
\midrule
\texttt{sab\_cache\_entries}     & $0$        & $1$        \\
\texttt{sab\_cache\_hits}        & $0$        & $1$        \\
\texttt{nuc\_cache\_entries}     & $0$        & $19$       \\
\texttt{nuc\_cache\_hits}        & $0$        & $19$       \\
\texttt{nuc\_cache\_bytes}       & $0$        & $2{,}468{,}534{,}728$\,B\\
\texttt{sab\_cache\_bytes}       & $0$        & $300$\,B   \\
\texttt{bundle\_cache\_budget}   & $3{,}220{,}881{,}408$\,B & (same)\\
\bottomrule
\end{tabular}}

\noindent
Reading: $2.47\,$GB of per-nuclide device buffers cached --- 19
distinct nuclides for one HEU-COMP-INTER case, under the
3.0\,GB budget. The 19 hits across one case correspond to one
hit per nuclide from the second seed; without the cache, the second
seed would have re-uploaded all 19 nuclide payloads. The
\texttt{sab\_cache} weighs only 300\,B for this case because the
case's TSL binding is narrow; a Godiva or pure-Be case with a
heavily-bound \texttt{c\_Be} TSL pushes that into the tens of MB.

\subsection{Why we did not yet add \texttt{material\_data} caching}

The material-composition upload
(\texttt{upload\_material\_data}: per-material nuclide indices,
atom densities, AWR table, $\bar{\nu}$ constants) is the third
production GPU upload. It also rebuilds on every call. We did not
cache it. The payload is in the low hundreds of KB for a
realistic ICSBEP case, roughly three orders of magnitude smaller
than the SAB payload it sits next to. The wasted bandwidth across
a sweep is dominated by the SAB payload by the same factor. We
chose to leave the engineering surface smaller until profile data
shows the material-data path is non-trivial.
